\chapter{Framework Architecture}
\label{ch:framework}

\section{Design Principles}

The framework is designed around four key principles:

\begin{enumerate}
    \item \textbf{Modularity}: Physics and ML components are decoupled, enabling independent development and ablation studies
    \item \textbf{Physical plausibility}: The model should respect known physical constraints (energy conservation, force balance)
    \item \textbf{Uncertainty quantification}: Probabilistic outputs with calibrated confidence intervals
    \item \textbf{Efficiency}: Inference time orders of magnitude faster than high-fidelity simulation
\end{enumerate}

% TODO: Add architecture diagram showing component interactions

\section{Baseline Physics Layer}

The physics layer implements a reduced-order model that approximates train dynamics using simplified force balance equations. For each car $i$:
\begin{equation}
    m_i \ddot{x}_i = F_{\text{trac},i} - F_{\text{res},i} - F_{\text{c},i-1} + F_{\text{c},i}
    \label{eq:physics_force_balance}
\end{equation}
where $F_{\text{trac},i}$ is traction/braking, $F_{\text{res},i}$ is resistance (Davis equation), and $F_{\text{c},i}$ is coupler force.

The physics layer serves two roles:
\begin{itemize}
    \item \textbf{Reference prediction}: Provides a baseline that the transformer can improve upon
    \item \textbf{Feature generation}: Computes intermediate quantities (forces, accelerations) that serve as input features to the transformer
\end{itemize}

% TODO: Detail the specific reduced-order model implementation
% TODO: Add table comparing physics layer accuracy vs full simulator

\section{Transformer Surrogate}

\subsection{Tokenization Scheme}

The transformer operates on a sequence of tokens representing:
\begin{itemize}
    \item \textbf{Route tokens}: One per segment, encoding grade, curvature, speed limit
    \item \textbf{Vehicle tokens}: One per car, encoding mass, length, position in consist
    \item \textbf{Control tokens}: One per time step, encoding traction/braking commands
    \item \textbf{Physics feature tokens}: Intermediate quantities from the physics layer
\end{itemize}

% TODO: Detail embedding dimensions and positional encoding

\subsection{Attention Factorization}

To manage computational complexity, we use factorized attention patterns:
\begin{itemize}
    \item \textbf{Temporal attention}: Models dependencies across time steps within the same car
    \item \textbf{Inter-car attention}: Models force propagation along the consist
    \item \textbf{Cross-attention}: Allows vehicles to attend to relevant route segments
\end{itemize}

This factorization reduces complexity from $O(n^2 m^2 T^2)$ to $O(n^2 + m^2 + T^2)$ where $n$ is segments, $m$ is cars, and $T$ is time steps.

% TODO: Add figure showing attention patterns
% TODO: Compare factorized vs full attention in ablation

\subsection{Output Heads}

The transformer decoder produces outputs through specialized heads:
\begin{itemize}
    \item \textbf{Quantile head}: Predicts $\quantile{0.05}, \quantile{0.25}, \quantile{0.50}, \quantile{0.75}, \quantile{0.95}$ for each output variable
    \item \textbf{Exceedance classifier}: Binary classifier for $\exceedprob{\couplerforce{i}(t)}{\tau}$ with threshold $\tau$
    \item \textbf{Deterministic head}: Point predictions for energy and arrival time
\end{itemize}

% TODO: Detail network architecture (layers, dimensions, activation functions)

\section{Hybrid Integration}

\subsection{Residual Learning}

The transformer predicts residuals relative to the physics layer:
\begin{equation}
    \hat{y} = y_{\text{phys}} + \Delta y_{\text{transformer}}
    \label{eq:residual_learning}
\end{equation}
This approach leverages the physics baseline while allowing the transformer to correct systematic errors.

\subsection{Physics Consistency Penalties}

To ensure physical plausibility, we add penalty terms:
\begin{itemize}
    \item \textbf{Energy conservation}: Predicted energy should match work done by forces
    \item \textbf{Force balance}: Sum of forces on each car should equal mass times acceleration
    \item \textbf{Bound constraints}: Coupler forces must be within physical limits
\end{itemize}

% TODO: Define penalty formulations mathematically
% TODO: Ablation study: with vs without physics penalties

\section{Inference Modes}

The framework supports two inference modes:

\begin{enumerate}
    \item \textbf{Single rollout}: Predict trajectory for one $(\route, \consist, \controlplan)$ tuple
    \item \textbf{Batch scenario sweeps}: Efficiently evaluate multiple scenarios for sensitivity analysis or optimization
\end{enumerate}

% TODO: Benchmark inference times for both modes
% TODO: Compare to high-fidelity simulation speedup
